\section{相关技术基础}

课程知识抽取的核心任务是从非结构化的课程文本中识别命名实体并抽取实体间语义关系，进而构建结构化的知识图谱。这一任务链条涉及多项关键技术：大语言模型提供基础的语义理解与文本生成能力；命名实体识别技术解决专业术语的定位与分类问题；关系抽取技术解决实体间语义关联的识别问题；知识图谱构建技术解决实体规范化、知识融合与大规模图谱的存储查询问题；教育领域知识图谱提供领域特定的构建方案与应用范式。本章按照"基础模型→核心抽取→图谱构建→领域应用"的技术脉络逐一展开，为第三、四章的系统设计与实现提供技术支撑。

\subsection{课程知识抽取技术总览}

\subsubsection{课程知识抽取面临的技术问题}

课程知识抽取不同于通用领域的信息抽取，其在技术层面面临一系列特有的困难。第一，领域实体边界模糊——课程术语常以复合形式出现（如"TCP拥塞控制""慢启动算法"），模型需要在专业知识与短语切分之间做出判断。第二，实体与关系类型多样——计算机网络课程涵盖协议、概念、机制、性能指标等不同知识元素，各类元素之间的语义关系类型差异较大。第三，跨文档实体统一——同一实体在不同教材和课件中以中英文全称、缩写等不同形式出现，需要有效的规范化机制。第四，长文本处理受限——完整教材往往包含数十万词元，而大语言模型的上下文窗口通常仅4K--8K词元，需要在分块策略与语义完整性之间取得平衡。

\subsubsection{技术路线图}

上述技术问题的解决需要多个技术层面的协同。图\ref{fig:2-1}给出了课程知识抽取的技术路线图，从底层大语言模型基础能力出发，经过命名实体识别和关系抽取两个核心技术层，最终汇聚到知识图谱构建与应用层。

\begin{figure}[ht]
    \centering
    \begin{tikzpicture}[
        node distance=0.4cm and 0.6cm,
        layer/.style={
            rectangle, rounded corners=5pt, draw=blue!60!black, fill=blue!6,
            minimum width=8.5cm, minimum height=1.8cm, align=center, font=\small
        },
        subbox/.style={
            rectangle, rounded corners=3pt, draw=#1!60!black, fill=#1!10,
            minimum height=0.6cm, minimum width=1.9cm, align=center,
            font=\footnotesize, text=black!80
        },
        arr/.style={-{stealth}, thick, draw=gray!60!black},
        dash/.style={thick, dashed, draw=gray!50!black, -{stealth}}
    ]
    % Layer 4: Education KG Application
    \node[layer] (l4) {
        \textbf{知识图谱构建与应用}\\[0.2cm]
        \begin{tikzpicture}
            \node[subbox=teal] (kgpipe) {流水线框架\\{\scriptsize InstructUIE, OneKE}};
            \node[subbox=teal, right=0.3cm of kgpipe] (kgcanon) {实体规范化\\{\scriptsize EDC, KGGen, SAC-KG}};
            \node[subbox=teal, right=0.3cm of kgcanon] (kgedu) {教育知识图谱\\{\scriptsize Edu-KG-Pipeline, Graphusion}};
        \end{tikzpicture}
    };

    % Layer 3: Relation Extraction
    \node[layer, below=0.3cm of l4] (l3) {
        \textbf{关系抽取技术}\\[0.2cm]
        \begin{tikzpicture}
            \node[subbox=orange] (re_few) {少样本RE\\{\scriptsize GPT-RE, Self-Prompting}};
            \node[subbox=orange, right=0.3cm of re_few] (re_doc) {文档级RE\\{\scriptsize AutoRE, Revisiting-RE}};
            \node[subbox=orange, right=0.3cm of re_doc] (re_zero) {零样本RE\\{\scriptsize GLiREL, REPaL}};
        \end{tikzpicture}
    };

    % Layer 2: Named Entity Recognition
    \node[layer, below=0.3cm of l3] (l2) {
        \textbf{命名实体识别技术}\\[0.2cm]
        \begin{tikzpicture}
            \node[subbox=violet] (ner_zero) {零样本NER\\{\scriptsize SLIMER, UniversalNER, GLiNER}};
            \node[subbox=violet, right=0.3cm of ner_zero] (ner_few) {少样本NER\\{\scriptsize GPT-NER, Self-Improving}};
            \node[subbox=violet, right=0.3cm of ner_few] (ner_code) {代码提示\\{\scriptsize CodeIE, GoLLIE, KnowCoder}};
        \end{tikzpicture}
    };

    % Layer 1: LLM Foundation
    \node[layer, below=0.3cm of l2] (l1) {
        \textbf{大语言模型基础}\\[0.2cm]
        \begin{tikzpicture}
            \node[subbox=red] (trans) {Transformer架构\\{\scriptsize 自注意力, 多头注意力}};
            \node[subbox=red, right=0.3cm of trans] (pretrain) {预训练与ICL\\{\scriptsize GPT系列, 上下文学习}};
            \node[subbox=red, right=0.3cm of pretrain] (prompt) {提示工程\\{\scriptsize Zero-Shot, Few-Shot, CoT}};
            \node[subbox=red, right=0.3cm of prompt] (mm) {多模态LLM\\{\scriptsize Qwen-VL, VLM OCR}};
        \end{tikzpicture}
    };

    % Arrows between layers
    \draw[arr] (l1.north) -- (l2.south);
    \draw[arr] (l2.north) -- (l3.south);
    \draw[arr] (l3.north) -- (l4.south);

    % Layer labels
    \node[left=0.5cm of l1, font=\footnotesize, text=gray!60, rotate=90] {基础层};
    \node[left=0.5cm of l2, font=\footnotesize, text=gray!60, rotate=90] {抽取层};
    \node[left=0.5cm of l3, font=\footnotesize, text=gray!60, rotate=90] {抽取层};
    \node[left=0.5cm of l4, font=\footnotesize, text=gray!60, rotate=90] {构建层};
    \end{tikzpicture}
    \caption{课程知识抽取技术路线图}\label{fig:2-1}
\end{figure}

如图\ref{fig:2-1}所示，课程知识抽取技术体系可划分为四个层次：（1）基础层——以大语言模型的Transformer架构、预训练范式、提示工程技术和多模态能力为底层支撑；（2）NER抽取层——通过零样本、少样本和代码提示三类范式解决课程术语的实体识别问题；（3）RE抽取层——通过少样本、文档级和零样本关系抽取技术识别实体间的语义关联；（4）构建层——通过流水线框架、实体规范化和教育领域知识图谱技术，将抽取结果整合为结构化的知识图谱。本章后续各节将按照这一技术路线图，从基础层出发，逐层展开各技术的原理、代表性方法和适用场景。

\subsection{大语言模型基础}

大语言模型（Large Language Models, LLMs）在超大规模文本语料上进行自监督预训练，具备语言理解与生成能力。这类模型有两个关键特性：涌现能力与上下文学习能力。本节从架构原理、预训练范式和多模态扩展三个层面展开。

\subsubsection{Transformer架构}

Transformer\cite{vaswani2017attention}用自注意力取代了循环结构，通过全并行计算绕过了RNN的串行瓶颈。其核心贡献有三点：自注意力机制（Self-Attention）、多头注意力（Multi-Head Attention）与位置编码（Positional Encoding）。

自注意力机制通过计算序列里头每个位置跟其他所有位置之间的注意力权重，让模型能够捕捉长距离的依赖关系。给定输入序列$\mathbf{X} = [\mathbf{x}_1, \mathbf{x}_2, ..., \mathbf{x}_n]$，自注意力的计算过程可以形式化为：
\begin{equation}
\text{Attention}(Q, K, V) = \text{softmax}\left(\frac{QK^T}{\sqrt{d_k}}\right)V
\end{equation}
其中$Q$、$K$、$V$分别代表查询（Query）、键（Key）和值（Value）矩阵，$d_k$为键向量的维度。缩放因子$\sqrt{d_k}$的用处在于缓解点积运算带来的梯度消失问题。自注意力机制可直接建模序列中任意两个位置之间的关系，不受距离限制，解决了RNN处理长序列时的梯度消失问题。

多头注意力机制将输入投影到多个子空间并行计算注意力，使模型从不同表示子空间中联合关注不同位置的信息。设有$h$个注意力头，那么多头注意力的输出就是：
\begin{equation}
\text{MultiHead}(Q, K, V) = \text{Concat}(\text{head}_1, ..., \text{head}_h)W^O
\end{equation}
其中每个$\text{head}_i = \text{Attention}(QW_i^Q, KW_i^K, VW_i^V)$，$W$为可学习的投影矩阵。多头机制让模型在不同子空间中关注不同类型的信息，如语法结构与语义关系，提升了模型的表达能力。不同注意力头确实学到了不同的语言特征：有的关注句法依赖，有的关注语义关联。

由于Transformer不含循环或卷积结构，引入位置编码以注入序列的顺序信息。原始Transformer采用正弦余弦函数来编码绝对位置：
\begin{align}
PE_{(pos, 2i)} &= \sin(pos/10000^{2i/d_{model}}) \\
PE_{(pos, 2i+1)} &= \cos(pos/10000^{2i/d_{model}})
\end{align}
后来的工作表明，可学习的绝对位置编码和相对位置编码在特定任务上表现更好。RoPE（Rotary Position Embedding）和ALiBi（Attention with Linear Biases）等相对位置编码方案在处理长序列时外推能力更强，可以支撑更长的上下文窗口。

\subsubsection{预训练语言模型与上下文学习}

在Transformer架构基础上，预训练语言模型形成了编码器和解码器两条主要技术路线。BERT采用双向Transformer编码器和掩码语言模型任务进行预训练，使模型能够同时利用左右上下文信息，推动了自然语言理解任务的发展\cite{devlin2019bert}。GPT（Generative Pre-trained Transformer）系列模型则采用自回归语言建模范式进行预训练。这种范式通过最大化序列的似然概率来学习语言表示，也就是：
\begin{equation}
\mathcal{L} = \sum_{i=1}^{n} \log P(x_i | x_1, ..., x_{i-1}; \theta)
\end{equation}
这种单向建模方式使得GPT系列模型特别擅长文本生成任务。预训练阶段，模型在大规模无标注语料上学习语言的统计规律和语义表示；下游任务中，再通过微调或提示工程将预训练知识迁移至特定应用场景。

GPT-3拥有1750亿参数，展示了上下文学习能力（In-Context Learning）\cite{brown2020gpt3}。上下文学习指模型不更新参数，仅凭输入中的少量示例即可适应新任务。这种能力源于预训练过程中对多样任务的隐式学习，模型通过示例模式泛化到待预测样本上。GPT-3在翻译、问答、摘要等任务上的性能接近监督学习方法，表明大模型有潜力作为通用任务求解器。

上下文学习的有效性高度依赖提示（Prompt）的设计。Brown等人发现，借助精心设计的任务描述和示例，GPT-3在多项自然语言理解任务上达到了接近微调模型的水平。这一发现催生了提示工程（Prompt Engineering）研究方向。后续工作分析了上下文学习的内在机制，包括示例选择策略、示例顺序的影响和标签空间覆盖的重要性等因素。

\subsubsection{提示工程技术}

提示工程是优化大语言模型输入以提升任务性能的一整套技术体系。按照所需示例的数量来分，提示策略可以分为零样本（Zero-Shot）、少样本（Few-Shot）和思维链（Chain-of-Thought）三种范式。

零样本提示向模型提供任务描述与待处理输入，无需任何标注示例，适用于快速原型验证和资源受限的场景。对于复杂任务，零样本提示的性能通常不足。指令微调（Instruction Tuning）通过在多样化的指令-响应对上微调模型，让模型更好地理解和执行新指令，从而改善了零样本性能。

少样本提示在输入中嵌入若干标注示例，引导模型理解任务模式。示例的选择、顺序和格式对性能影响很大。有效的少样本提示应确保示例与测试样本在分布上保持一致，并覆盖任务的主要变化模式。KATE、EPR等方法通过检索与测试样本相似的示例来优化提示；Self-Generated In-Context Learning则让模型自己生成示例，降低对人工标注的依赖。

思维链提示由Wei等人提出，通过在示例中展示中间推理步骤，激发模型的逐步推理能力。该方法将复杂问题拆分为一系列可解释的推理步骤，提升了模型在算术推理、常识推理和符号推理任务上的表现。后续工作在此基础上发展出自一致性解码（Self-Consistency）、思维树（Tree of Thoughts）等增强技术。自一致性通过多次采样并选出最一致的答案来提升推理可靠性；思维树则通过探索多条推理路径并评估各路径的可行性，实现更复杂的决策过程。

\subsubsection{多模态大语言模型}

多模态大语言模型（Multimodal LLMs）将纯文本模型的能力扩展到视觉、音频等多种模态，使模型能够理解并生成多模态内容。在知识抽取领域，视觉-语言模型（Vision-Language Models）尤其重要，因为它可以直接处理扫描文档、幻灯片等图像格式的课程材料，无需额外的OCR预处理步骤。

典型的视觉-语言模型架构包含视觉编码器、投影层和语言模型三个部分。视觉编码器（如ViT、CLIP）将输入图像编码为视觉特征序列；投影层负责将视觉特征映射到语言模型的语义空间；语言模型基于融合后的多模态表示生成文本输出。早期视觉-语言模型如BLIP、Flamingo等通过端到端训练实现跨模态对齐；LLaVA、Qwen-VL等工作采用两阶段训练策略：先进行视觉-语言对齐预训练，再进行指令微调，使模型能够遵循指令。

Qwen-VL、InternVL等模型在文档理解、图表解析等任务上表现好，可直接用于教育材料自动化处理。这些模型能识别图像中的文本内容、理解版面布局、解析表格和图示结构，提升了从非结构化文档中提取结构化信息的能力。对于包含复杂排版和数学公式的学术材料，多模态模型的端到端处理能力优于传统的OCR加文本处理流水线。

\subsection{命名实体识别技术}

命名实体识别（Named Entity Recognition, NER）的目标是从非结构化文本中识别命名实体并对其分类。传统方法依赖大量标注数据和特定领域的特征工程，而基于大模型的方法借助提示工程实现了低资源、甚至零资源场景下的实体识别。按照利用训练数据的方式，现有方法分为零样本、少样本和代码提示三类范式\cite{chewanxiang2021}。

\subsubsection{零样本命名实体识别}

零样本NER指在没有任何标注样本的情况下，仅靠任务描述和实体类型定义完成实体识别，适用于快速适应新领域、新实体类型。SLIMER\cite{slimer_2024}用丰富的实体定义和标注指南弥补训练数据的缺失，采用单实体类型每查询（Single Entity Type Per Query）策略，将多类型NER任务拆分为多个二分类子任务，降低了任务复杂度。该方法通过详细的类型描述、边界说明和示例指导模型做精确识别，在多个基准数据集上比传统零样本方法更好。

UniversalNER探索了运用ChatGPT进行自动标注并蒸馏到小模型的技术路线\cite{universalner_2023}。该方法先用ChatGPT生成大规模弱监督数据，再通过知识蒸馏将大模型的能力迁移到轻量级学生模型上。具体来说，UniversalNER设计了一系列针对NER任务的提示模板，引导ChatGPT生成高质量的伪标注数据，然后使用这些数据训练专门的NER模型。实验表明，这种范式在开放域NER任务上有竞争力，同时降低了推理成本。

GLiNER提出了一种基于双向Transformer的实体-跨度匹配架构\cite{gliner_2023}。与传统序列标注方法不同，GLiNER将NER建模为实体类型与文本跨度之间的匹配问题，通过计算类型表示和跨度表示的相似度预测实体边界和类型。GLiNER将实体类型视为可学习的嵌入，使其泛化到训练时未见过的全新类型，在零样本和少样本场景下均有效，支持任意实体类型的灵活扩展。

\subsubsection{少样本命名实体识别}

少样本NER利用少量标注示例引导大模型适应特定领域和实体类型。GPT-NER\cite{gpt_ner_2023}使用\texttt{@@}和\texttt{\#\#}标记实体边界，将NER转化为大模型擅长的文本生成任务。这种标记方式与GPT模型的预训练目标更契合，能更好地激发模型的标注能力。GPT-NER还引入了三层次演示检索策略，分别从句子级、实体级和词级检索相似示例，确保示例在多个粒度上与测试样本相关。此外，GPT-NER通过自验证机制过滤低质量预测结果，提高了输出的可靠性。

Self-Improving NER\cite{self_improving_ner_2023}通过自一致性投票和可靠性筛选实现模型性能的自我迭代提升。该方法先用不同提示变体生成多组预测结果，通过一致性投票找出高置信度预测；再基于可靠性指标筛选优质伪标签用于下一轮训练。这种自举式学习策略降低了对人工标注的依赖，使模型从自身的预测中持续学习和改进。经过若干轮迭代后，模型性能持续提升，在某些任务上接近全监督方法的表现。

LLM-DA（LLM-based Data Augmentation）从上下文和实体两个层面进行数据增强\cite{llm_da_2024}。上下文层面运用大模型生成与原始句子语义相似的变体，同时保持实体位置不变；实体层面通过实体替换和同义词扩展来丰富实体表达的多样性。增强后的数据与原始少样本数据联合训练，提升了模型的泛化能力。LLM-DA借助大模型的语言理解能力进行语义保持的数据增强，生成的样本质量优于传统基于规则的增强方法。

\subsubsection{代码提示范式}

Code-as-Prompt是一类较新的提示范式，将信息抽取任务表示为代码形式，利用大模型的代码理解能力提升任务性能。CodeIE\cite{codeie_2023}将NER和关系抽取任务转化为Python代码生成任务。实体识别被建模为从文本中提取特定类型值的函数实现，关系抽取被建模为返回实体对的函数。实验表明，代码形式的提示在自然语言理解任务上优于纯文本提示，这可能与代码的结构化特性以及大模型在代码预训练中获得的能力有关。代码表示能更清晰地表达任务的结构和约束条件，减少了自然语言歧义带来的理解偏差。

GoLLIE扩展了代码提示范式\cite{gollie_2023}，引入Python类定义来表示实体和关系的模式（Schema）。每个实体类型和关系类型被定义为具有特定属性和文档字符串的Python类，文档字符串中详细描述了类型的定义和标注指南。这种表示方式将模式定义与提示模板统一起来，支持复杂嵌套结构的抽取。GoLLIE在多个IE基准测试上达到了当时的SOTA。KnowCoder\cite{knowcoder_2024}在此基础上构建了大规模代码风格知识模式库，通过代码预训练与指令微调两阶段学习，将通用预训练的知识抽取能力迁移到特定领域。

\subsection{关系抽取技术}

关系抽取（Relation Extraction, RE）的目标是识别文本中实体之间的语义关系。与命名实体识别相比，关系抽取需要建模实体对之间的复杂交互，任务难度更高。基于大模型的关系抽取方法在少样本和零样本场景下有优势，可借助预训练知识快速适应新的关系类型和领域。

\subsubsection{少样本关系抽取}

GPT-RE针对少样本关系抽取提出了任务感知演示检索和金标签诱导推理两种技术\cite{gpt_re_2023}。任务感知检索基于实体对和关系标签的联合相似度选择演示示例，确保示例与查询在实体分布和关系模式上高度一致；金标签诱导推理在提示中显式引导模型关注实体提及的位置，并通过标签的语义信息约束预测空间。两种技术的结合使GPT-RE在FewRel等基准上性能提升明显。GPT-RE还探讨了不同检索策略和提示格式对性能的影响。

Self-Prompting\cite{self_prompting_re_2024}通过三阶段流程实现自动化的少样本关系抽取。第一阶段运用大模型生成多样化的合成样本，覆盖关系类型的多种表达模式；第二阶段对生成的样本进行质量筛选和去重，保留高质量的伪标注数据；第三阶段基于筛选后的样本自动构建提示模板，无需人工设计。这种自引导方法降低了对人工编写示例的依赖。Self-Prompting的核心洞察在于，大模型本身具备生成合理训练样本的能力，通过适当的引导和筛选，可构建有效的自我监督流程。

\subsubsection{文档级关系抽取}

文档级关系抽取需要处理跨越多个句子的长文本，识别显式提到或隐式蕴含的实体关系。这项任务比句子级关系抽取更具挑战性，需要建模复杂的指代消解和跨句推理能力。AutoRE提出了关系-头部-事实（Relation-Head-Fact, RHF）的统一范式\cite{auto_re_2024}，将关系抽取分解为关系识别、头部实体定位和事实抽取三个子任务。该方法采用QLoRA进行参数高效微调，在DocRED等文档级基准上性能好。RHF范式将复杂的文档级关系抽取分解为多个更易管理的子任务，降低了学习难度，同时提升了可解释性。

Revisiting-RE系统分析了文档级关系抽取中的关键挑战\cite{revisiting_re_2023}，提出将思维链增强和人工评估相结合的混合策略。直接使用大模型进行文档级关系抽取会面临实体指代消解困难和长距离依赖建模复杂等问题；引入逐步推理提示和人工反馈的迭代优化可缓解这些难题。该工作特别在评估协议和错误分析方面做了深入探讨，指出了当前方法的局限性与改进方向。

\subsubsection{零样本关系抽取}

GLiREL延续了GLiNER的设计理念，将零样本关系抽取建模为关系类型与实体对之间的匹配问题\cite{glirel_2025}。该模型选用轻量级双向Transformer架构，单次前向传播即可计算所有候选实体对和关系类型的匹配分数，实现了高效推理。GLiREL支持任意关系类型的灵活扩展，无需针对新类型重新训练，实用性高。该方法将关系类型编码为可学习的嵌入向量，通过计算实体对表示与关系类型表示的兼容性来预测关系。

REPaL（Relation Extraction with Pairwise Loss）探索了仅靠关系定义完成零样本抽取的可能性\cite{repal_2024}。该方法运用大模型生成关系定义的语义嵌入，通过对比学习使相关实体对的表示靠近、无关实体对的表示远离。REPaL在零样本文档级关系抽取任务上性能领先，证明关系定义本身蕴含丰富的语义信息。REPaL完全不依赖标注样本，仅凭关系类型的自然语言描述即可进行抽取，适合快速搭建原型系统以及处理新增的关系类型。

\subsection{知识图谱构建与实体规范化技术}

知识图谱构建涉及实体识别、关系抽取、实体链接、知识融合等多个环节。近年来，大语言模型凭借其语义理解能力，在知识图谱构建各环节中得到了应用。本节从流水线框架和实体规范化两个方向介绍相关技术进展\cite{zhaojun2018,wanghaofen2019}。

\subsubsection{流水线框架}

LLMs-KG-Survey系统综述了运用大模型构建知识图谱的研究进展\cite{llms_kg_survey_2023}。该综述将现有方法分为自主构建和协同构建两类，分析了各类方法的优势和局限。此外，AutoKG多智能体框架通过规划、抽取、验证、修正等多个智能体之间的协作，实现了端到端的知识图谱构建。AutoKG采用模块化设计，各智能体专注于特定子任务并通过消息传递进行协作，提升了系统的可扩展性和鲁棒性。

InstructUIE采用多任务指令微调策略统一处理信息抽取任务\cite{instructuie_2023}。该方法将NER、关系抽取、事件抽取等多种任务统一为文本到文本的生成任务，借助自然语言指令描述任务需求。指令微调使模型具备任务泛化能力，在未见过的任务类型和领域上仍能维持合理性能。

OneKE提出了基于Schema引导的多智能体知识抽取框架\cite{oneke_2025}。该框架先用大模型从领域文档中自动抽取Schema定义，再基于Schema指导抽取智能体识别实体和关系，最后通过验证智能体保证抽取结果的一致性和完整性。OneKE在教育、医疗等垂直领域适应性好。RUIE\cite{ruie_2024}从检索增强的角度探索了统一信息抽取，通过双编码器检索器配合LLM偏好评分实现跨任务的演示检索，在低资源场景下降低了模型微调成本。

\subsubsection{实体规范化}

实体规范化（Entity Canonicalization）的目标是识别指代同一现实实体的不同提及方式，是知识图谱构建的重要环节。EDC提出了抽取-定义-规范化的三阶段方法\cite{edc_2024}：先从文本中抽取候选实体，再运用大模型为每个实体生成语义定义，最后基于定义相似度进行实体聚类和规范化。该方法处理了实体表达的多样性问题。通过引入语义定义作为中间表示，EDC能捕捉实体提及的深层语义信息，而不只是依赖表面的字符串相似度。

KGGen探索了聚类和大模型混合迭代的实体去重策略\cite{kggen_2025}，先用传统聚类算法进行初步分组，再运用大模型对每个聚类进行语义审核和合并决策，迭代直至收敛。这种混合策略兼顾了效率和准确性，在大规模知识图谱构建中效果好。KGGen将传统聚类的高效性与大模型的语义理解能力相结合，通过迭代优化提升了规范化质量。

SAC-KG设计了生成器-验证器-剪枝器的三组件架构\cite{sac_kg_2024}。生成器负责从源文本中抽取候选实体和关系；验证器基于大模型评估三元组的事实性和一致性；剪枝器移除低质量或重复的三元组。三个组件协同工作，确保知识图谱的高精度和低冗余。LOKE\cite{loke_2023}从开放知识抽取与实体链接的角度提出了另一种方案，通过编辑距离算法将抽取的三元组链接到Wikidata，在保障三元组完整性的同时实现了实体可链接性。

\subsection{教育领域知识图谱}

教育领域知识图谱面向课程教材、学术论文等教育资源，旨在构建结构化的学科知识体系，支撑智能问答、学习路径推荐等教育应用。与通用领域知识图谱相比，教育领域知识图谱的概念层次更清晰、知识关联更密集、应用场景更特定。本节梳理近年来教育领域知识图谱构建的代表性工作。

国内研究也开始从智能教育角度讨论知识图谱与大语言模型的协同共生关系，认为知识图谱可为大模型提供结构化知识约束，而大模型可降低知识图谱构建与维护成本\cite{lixiaoli2025_kg_llm_education_survey}。这与课题面向课程材料自动抽取知识并进行证据约束的技术路线一致。

Educational-KG-Pipeline提出了面向教育资源的五阶段知识图谱构建流水线\cite{educational_kg_pipeline_2025}，包括文档解析、实体识别、关系抽取、实体链接和知识融合，针对PDF、PPT等格式的课程材料进行了专门优化。文档解析阶段处理了教育材料中常见的复杂版面和多媒体元素；后续阶段针对教育领域的实体类型和关系模式设计了专门的抽取策略。在计算机网络课程数据集上，该流水线将F1分数从40\%提升至47\%。

Graphusion专注于概念知识图谱的构建\cite{graphusion_2024}，采用种子概念扩展、候选三元组生成、全局融合的三阶段策略。该方法从少量人工标注的种子概念出发，运用大模型自动发现相关概念和关系；借助全局一致性约束解决冲突和冗余问题；最后构建了TutorQA基准数据集用于教育问答评估。Graphusion在概念覆盖度和关系准确性方面表现好。该工作通过种子引导和自动扩展减少了对大规模人工标注的依赖，是一种低成本的概念图谱构建方法。

LLM-Curriculum-KG探索了人机协作的课程知识图谱构建模式\cite{llm_curriculum_kg_2025}，设计了课程-领域-用户三层本体结构：课程层描述课程大纲和章节组织，领域层涵盖学科概念和知识关联，用户层记录学习者的知识状态和学习轨迹。大模型负责自动抽取和初步验证，人工专家进行质量审核和领域知识补充。这种人机协作模式兼顾了大模型的自动化处理能力和人类专家的专业判断。

\subsection{本章小结}

本章梳理了课程知识命名实体识别与关系抽取的关键技术。首先给出了课程知识抽取的技术总览，建立了"基础层→NER层→RE层→构建层"的四层技术路线图。随后，大语言模型部分聚焦Transformer的自注意力机制、GPT系列的自回归预训练范式，以及多模态视觉-语言模型；命名实体识别部分对比了零样本、少样本和代码提示三类范式的核心思路与适用场景；关系抽取部分讨论了句子级与文档级场景下的任务感知检索、自引导提示与零样本匹配等方法；知识图谱构建部分介绍了流水线框架与实体规范化技术；教育领域部分分析了现有知识图谱方案的特点与局限。
